\section{Zasilanie i bilans mocy}
\lead{Terminal ma jeden centralny UPS dla rdzenia sieci, zapisu i SOC. Szafy polowe SP1 i SP2 mają własne obwody 230 V z rozdzielnicy przy N0 i lokalne UPS na szynie DIN. Rozdział podaje moc do pokrycia, a czas podtrzymania liczy rozdział 13.}

\fig{fig-chart-power}{0.9}{Bilans mocy: obwiednia AC centralnego UPS oraz obciążenie zasilaczy 48 V w szafach polowych z zapasem 25\%.}

\subsection*{Centralny UPS}
Do centralnego UPS należą tylko urządzenia w pomieszczeniu IT i w SOC. Kamery terminalu, KP01, kontrolery KD z czytnikami i zaczepami są odbiorami PoE przełącznika N0, więc ich moc zawiera się w poborze AC przełącznika.
\begin{center}
\small
\begin{tabularx}{\linewidth}{@{}lrY@{}}
\toprule
Odbiornik & AC [W] & Podstawa \\
\midrule
N0, TP-Link SG3428MP & 478,6 & najgorszy pobór AC z pełnym budżetem PoE 384 W \\
NAS, Synology DS923+ & 125,0 & zasilacz 100 W, przyjęta sprawność 0,8 \\
Komputer SOC z serwerem i klientem VMS & 112,5 & zasilacz 90 W, przyjęta sprawność 0,8 \\
Dwa monitory, Dell P2425H & 130,0 & 65 W maksymalnie na monitor \\
Zapora, TP-Link ER605 & 15,0 & 7,94 W urządzenia, granica z zasilaczem \\
\midrule
Razem & 861,1 & z zapasem 25\%: 1076 W \\
\bottomrule
\end{tabularx}
\end{center}
Obwiednia N0 zakłada pełny budżet PoE 384 W, a suma mocy maksymalnych odbiorników PoE z rozdziału 7 wynosi 252 W. Różnica 132 W zostaje jako rezerwa na dodatkowe urządzenia PoE. Sprawności zasilaczy NAS, komputera i zapory są granicami przyjętymi w projekcie, a nie wartościami z kart. Monitory policzono mocą maksymalną, a nie typową.\cite{S-SW,S-NAS,S-PC,S-MON,S-FW}

UPS dobrany w rozdziale 13 musi mieć moc wyjściową co najmniej 1076 W. Nie ma osobnego serwera analityki, bo analityka działa w kamerach (rozdział 6). SSP i SSWiN mają własne zasilacze z akumulatorami i nie są odbiorami UPS.\cite{S-FIRE,S-SSWIN}

\subsection*{Szafy polowe}
Przełącznik T8504-R pracuje z napięcia od 48 do 57 V DC. Zasila go zasilacz DIN PS56 240 W, a zasilacz pracuje z lokalnego UPS na szynie DIN.\cite{S-T8504} Karta przełącznika nie podaje poboru własnego, więc przyjęto 20 W. Odbiory PoE liczy się według przydziału portów z rozdziału 7.\cite{OWN}
\begin{center}
\small
\begin{tabular}{@{}lrr@{}}
\toprule
Pozycja & SP1 [W] & SP2 [W] \\
\midrule
Przydział PoE & 120,0 & 15,4 \\
Przełącznik, wartość przyjęta & 20,0 & 20,0 \\
\midrule
Obciążenie zasilacza 48 V & 140,0 & 35,4 \\
Z zapasem 25\% & 175,0 & 44,3 \\
Pobór AC zasilacza, sprawność 0,9 & 155,6 & 39,3 \\
Wyjście UPS DIN z zapasem 25\% & 194,4 & 49,2 \\
\bottomrule
\end{tabular}
\end{center}
Zasilacz SP1 jest obciążony w 73\% z zapasem, a SP2 w 18\%. Obie szafy mogą mieć ten sam UPS DIN o mocy wyjściowej co najmniej 195 W, co upraszcza serwis.

Napęd bramy B1 to przykładowo FAAC 741: 230 V, 500 W, 2,2 A, skrzydło do 900 kg.\cite{S-GATE} Napęd ma osobny wyłącznik w SP1 przed UPS DIN. Przy zaniku sieci brama pozostaje zamknięta, a patrol może ją otworzyć ręcznie po odblokowaniu napędu. Furtka z czytnikiem działa dalej, bo kontroler B1 jest zasilany z PoE.\cite{S-KD}

\subsection*{Obwody 230 V do szaf}
Każda szafa ma osobny obwód z rozdzielnicy przy N0, prowadzony trasą T3 (rozdział 9). Prąd obwodu z zapasem 25\% obejmuje też ładowanie UPS DIN:
\[
I_{\mathrm{SP1}} = 1{,}25\left(\frac{155{,}6}{230\cdot0{,}9}+2{,}2\right)=3{,}7\ \mathrm{A},\qquad
I_{\mathrm{SP2}} = 1{,}25\cdot\frac{39{,}3}{230\cdot0{,}9}=0{,}24\ \mathrm{A}.
\]
Spadek napięcia dla obwodu jednofazowego, z $\gamma=56\ \mathrm{m}/(\Omega\,\mathrm{mm}^2)$ dla miedzi i $\cos\varphi=1$ jako przypadkiem niekorzystnym:
\[
\Delta U\% = \frac{200\,L\,I}{\gamma\,S\,U}.
\]
\begin{center}
\small
\begin{tabular}{@{}lrrrr@{}}
\toprule
Obwód & $L$ [m] & $I$ [A] & $S$ [mm$^2$] & $\Delta U$ [\%] \\
\midrule
SP1 & 166,6 & 3,7 & 2,5 & 3,8 \\
SP1, przyjęty & 166,6 & 3,7 & 4 & 2,4 \\
SP2, przyjęty & 317,6 & 0,24 & 2,5 & 0,5 \\
\bottomrule
\end{tabular}
\end{center}
Przyjęto granicę 3\%.\cite{OWN} Dla SP1 przekrój 2,5 mm$^2$ jej nie spełnia, więc przyjęto kabel ziemny 3 $\times$ 4 mm$^2$. Dla SP2 wystarcza 3 $\times$ 2,5 mm$^2$. Prądu rozruchu napędu nie sprawdzano. Oba obwody mają wyłącznik różnicowoprądowy 30 mA, który zapewnia samoczynne wyłączenie przy długim przewodzie. Rzeczywiste impedancje pętli zwarcia podlegają pomiarom przed uruchomieniem.

\subsection*{Wyłącznik pożarowy i systemy z własnym zasilaniem}
Przeciwpożarowy wyłącznik prądu odłącza wejście AC centralnego UPS i otwiera jego wejście RPO osobnym bezpotencjałowym stykiem. Wejście RPO trzeba włączyć w konfiguracji UPS, bo fabrycznie nie ma przypisanej funkcji. Zwykły zanik sieci nie otwiera RPO, więc UPS podtrzymuje zasilanie.\cite{S-UPS-GUIDE} Obwody SP1 i SP2 wychodzą z rozdzielnicy za wyłącznikiem pożarowym. UPS DIN w szafach zasilają tylko urządzenia w szafie i przy OG1, a nie instalację budynku. SSP pracuje z własnego obwodu przed wyłącznikiem pożarowym i z własnej baterii.\cite{S-FIRE}
